\documentclass[10pt,a4paper]{amsart}
\usepackage{amsmath,amssymb,amsthm,graphicx}
\usepackage{verbatim}
\newenvironment{reference}{\comment}{\endcomment}
%\newenvironment{reference}{}{}
\newenvironment{slogan}{\comment}{\endcomment}
\newenvironment{history}{\comment}{\endcomment}

% For commutative diagrams we use Xy-pic
\usepackage[all]{xy}

% We use 2cell for 2-commutative diagrams.
\xyoption{2cell}
% Theorem environments.
%
\theoremstyle{plain}
\newtheorem{theorem}[subsection]{Theorem}
\newtheorem{proposition}[subsection]{Proposition}
\newtheorem{lemma}[subsection]{Lemma}

\theoremstyle{definition}
\newtheorem{definition}[subsection]{Definition}
\newtheorem{example}[subsection]{Example}
\newtheorem{exercise}[subsection]{Exercise}
\newtheorem{situation}[subsection]{Situation}

\theoremstyle{remark}
\newtheorem{remark}[subsection]{Remark}
\newtheorem{remarks}[subsection]{Remarks}

\numberwithin{equation}{subsection}

% Macros
%
\def\lim{\mathop{\mathrm{lim}}\nolimits}
\def\colim{\mathop{\mathrm{colim}}\nolimits}
\def\Spec{\mathop{\mathrm{Spec}}}
\def\Hom{\mathop{\mathrm{Hom}}\nolimits}
\def\Ext{\mathop{\mathrm{Ext}}\nolimits}
\def\SheafHom{\mathop{\mathcal{H}\!\mathit{om}}\nolimits}
\def\SheafExt{\mathop{\mathcal{E}\!\mathit{xt}}\nolimits}
\def\Sch{\mathit{Sch}}
\def\Mor{\mathop{\mathrm{Mor}}\nolimits}
\def\Ob{\mathop{\mathrm{Ob}}\nolimits}
\def\Sh{\mathop{\mathit{Sh}}\nolimits}
\def\NL{\mathop{N\!L}\nolimits}
\def\CH{\mathop{\mathrm{CH}}\nolimits}
\def\proetale{{pro\text{-}\acute{e}tale}}
\def\etale{{\acute{e}tale}}
\def\QCoh{\mathit{QCoh}}
\def\Ker{\mathop{\mathrm{Ker}}}
\def\Im{\mathop{\mathrm{Im}}}
\def\Coker{\mathop{\mathrm{Coker}}}
\def\Coim{\mathop{\mathrm{Coim}}}

% Boxtimes
%
\DeclareMathSymbol{\boxtimes}{\mathbin}{AMSa}{"02}

%
% Macros for moduli stacks/spaces
%
\def\QCohstack{\mathcal{QC}\!\mathit{oh}}
\def\Cohstack{\mathcal{C}\!\mathit{oh}}
\def\Spacesstack{\mathcal{S}\!\mathit{paces}}
\def\Quotfunctor{\mathrm{Quot}}
\def\Hilbfunctor{\mathrm{Hilb}}
\def\Curvesstack{\mathcal{C}\!\mathit{urves}}
\def\Polarizedstack{\mathcal{P}\!\mathit{olarized}}
\def\Complexesstack{\mathcal{C}\!\mathit{omplexes}}
% \Pic is the operator that assigns to X its picard group, usage \Pic(X)
% \Picardstack_{X/B} denotes the Picard stack of X over B
% \Picardfunctor_{X/B} denotes the Picard functor of X over B
\def\Pic{\mathop{\mathrm{Pic}}\nolimits}
\def\Picardstack{\mathcal{P}\!\mathit{ic}}
\def\Picardfunctor{\mathrm{Pic}}
\def\Deformationcategory{\mathcal{D}\!\mathit{ef}}
\usepackage{cleveref}
\crefname{originalreference}{original reference}{original references}
\Crefname{originalreference}{Original reference}{Original references}
\makeatletter
\newlabel{situation-dualizing}{{0AU4}{source}{}{sourcepointer-1}{}}
\newlabel{situation-dualizing@cref}{{[originalreference][0][]0AU4}{[1][0][]source}{}{}{}}
\makeatother
\crefname{originallemma}{original lemma}{original lemmas}
\Crefname{originallemma}{Original lemma}{Original lemmas}
\makeatletter
\newlabel{lemma-dualizing-schemes}{{0A89}{source}{}{sourcepointer-2}{}}
\newlabel{lemma-dualizing-schemes@cref}{{[originallemma][0][]0A89}{[1][0][]source}{}{}{}}
\makeatother
\crefname{originallemma}{original lemma}{original lemmas}
\Crefname{originallemma}{Original lemma}{Original lemmas}
\makeatletter
\newlabel{categories-lemma-properties-2-cat-cats}{{003F}{source}{}{sourcepointer-3}{}}
\newlabel{categories-lemma-properties-2-cat-cats@cref}{{[originallemma][0][]003F}{[1][0][]source}{}{}{}}
\makeatother
\crefname{originallemma}{original lemma}{original lemmas}
\Crefname{originallemma}{Original lemma}{Original lemmas}
\makeatletter
\newlabel{lemma-twisted-inverse-image}{{0A9E}{source}{}{sourcepointer-4}{}}
\newlabel{lemma-twisted-inverse-image@cref}{{[originallemma][0][]0A9E}{[1][0][]source}{}{}{}}
\makeatother
\crefname{originallemma}{original lemma}{original lemmas}
\Crefname{originallemma}{Original lemma}{Original lemmas}
\makeatletter
\newlabel{lemma-shriek-via-duality}{{0AU2}{source}{}{sourcepointer-5}{}}
\newlabel{lemma-shriek-via-duality@cref}{{[originallemma][0][]0AU2}{[1][0][]source}{}{}{}}
\makeatother
\crefname{originallemma}{original lemma}{original lemmas}
\Crefname{originallemma}{Original lemma}{Original lemmas}
\makeatletter
\newlabel{lemma-proper-map-good-dualizing-complex}{{0AUD}{source}{}{sourcepointer-6}{}}
\newlabel{lemma-proper-map-good-dualizing-complex@cref}{{[originallemma][0][]0AUD}{[1][0][]source}{}{}{}}
\makeatother
\crefname{originalsection}{original section}{original sections}
\Crefname{originalsection}{Original section}{Original sections}
\makeatletter
\newlabel{example-iso-on-RSheafHom-noetherian}{{0GEW}{source}{}{sourcepointer-7}{}}
\newlabel{example-iso-on-RSheafHom-noetherian@cref}{{[originalsection][0][]0GEW}{[1][0][]source}{}{}{}}
\makeatother
\title{The newly defined shriek functors form a pseudo functor and are compatible with the duality involution and with the right adjoint for a proper morphism}
\author{The Stacks Project Authors}
\date{stacks-project:duality.tex:local-snapshot-20261012; GNU FDL 1.2 or later}
\begin{document}
\maketitle
\noindent\textit{Editorial note.} All delivered mathematical text is the author's own, copied from the pinned Stacks Project chapter file duality.tex; no statement or proof step was rewritten, reordered or removed.\par\medskip
\noindent\textit{Editorial note.} Preamble provenance: the pinned chapter file opens with an input of the shared preamble and carries no preamble text of its own, so the project's single Stacks-wide preamble staged at sources/stacks-duality/preamble.tex is excerpted verbatim under the recipe fixed by LANE\_BRIEF 22.1, as three spans, lines 8-13 (verbatim package and the reference/slogan/history environments), lines 14-18 (Xy-pic setup) and lines 154-217 (the native theorem declarations and mathematical macros). Nothing was added, removed or redefined; the wrapper class and the amsmath/amssymb/amsthm/graphicx base packages are supplied by the assembler.\par\medskip
\noindent\textit{Editorial note.} The counted statement and proof contain no cross reference, no citation and no included file, so the delivered body is self contained.\par\medskip
\noindent\textit{Editorial note.} Canonical location of the counted proposition: Stacks tag 0AUE (sources/stacks-duality/pinned-original-source.tex line 5588).\par\medskip
\noindent\textit{Allowed standard prerequisites (editorial).}\begin{itemize}
\item The duality involution on the bounded coherent derived category. Same chapter, Lemma labelled lemma-dualizing-schemes, cited at proof lines 5640-5644 and 5668. Supplies eta\_Y, eta\_Z and part (3) of the statement.
\item The upper shriek through duality and the right adjoint. Same chapter, Lemma labelled lemma-shriek-via-duality, cited at proof lines 5674-5676, 5678-5680 and 5684-5686. Identifies f\textasciicircum{}!\_new with the right adjoint for the proper case and with f\textasciicircum{}! for the separated case.
\item Identification of the good dualizing complex with the shriek of the base complex for a proper morphism. Same chapter, Lemma labelled lemma-proper-map-good-dualizing-complex, cited at proof lines 5670-5674. Gives omega\_X\textasciicircum{}bullet = a(omega\_Y\textasciicircum{}bullet) needed to change the base.
\item The Noetherian example comparing the two internal Hom expressions. Same chapter, Example labelled example-iso-on-RSheafHom-noetherian, cited at proof lines 5682-5684. Supplies the displayed equalities of part (6).
\end{itemize}\medskip
\ifdefined\hypertarget\hypertarget{sourcepointer-1}{}\fi
\noindent\textit{Source reference.} Original reference 0AU4: sources/stacks-duality/pinned-original-source.tex, source lines 5259--5259. Same-chapter pointer into the author environment situation at source line 5259; invoked at proof line None of the counted proof. It is retained only as a navigation pointer; every load-bearing step is inside the delivered span.\par\medskip
\ifdefined\hypertarget\hypertarget{sourcepointer-2}{}\fi
\noindent\textit{Source reference.} Original lemma 0A89: sources/stacks-duality/pinned-original-source.tex, source lines 472--472. Same-chapter pointer into the author environment lemma at source line 472; invoked at proof line 5649 of the counted proof. It is retained only as a navigation pointer; every load-bearing step is inside the delivered span.\par\medskip
\ifdefined\hypertarget\hypertarget{sourcepointer-3}{}\fi
\noindent\textit{Source reference.} Original lemma 003F: sources/stacks-chow/tags.txt, source lines 132--132. Cross-chapter standard prerequisite located in the project tag list staged in the local snapshot; invoked at proof line 5650 of the counted proof. It is retained only as a navigation pointer; every load-bearing step is inside the delivered span.\par\medskip
\ifdefined\hypertarget\hypertarget{sourcepointer-4}{}\fi
\noindent\textit{Source reference.} Original lemma 0A9E: sources/stacks-duality/pinned-original-source.tex, source lines 653--653. Same-chapter pointer into the author environment lemma at source line 653; invoked at proof line 5668 of the counted proof. It is retained only as a navigation pointer; every load-bearing step is inside the delivered span.\par\medskip
\ifdefined\hypertarget\hypertarget{sourcepointer-5}{}\fi
\noindent\textit{Source reference.} Original lemma 0AU2: sources/stacks-duality/pinned-original-source.tex, source lines 4734--4734. Same-chapter pointer into the author environment lemma at source line 4734; invoked at proof line 5673 of the counted proof. It is retained only as a navigation pointer; every load-bearing step is inside the delivered span.\par\medskip
\ifdefined\hypertarget\hypertarget{sourcepointer-6}{}\fi
\noindent\textit{Source reference.} Original lemma 0AUD: sources/stacks-duality/pinned-original-source.tex, source lines 5541--5541. Same-chapter pointer into the author environment lemma at source line 5541; invoked at proof line 5676 of the counted proof. It is retained only as a navigation pointer; every load-bearing step is inside the delivered span.\par\medskip
\ifdefined\hypertarget\hypertarget{sourcepointer-7}{}\fi
\noindent\textit{Source reference.} Original section 0GEW: sources/stacks-duality/pinned-original-source.tex, source lines 865--865. Same-chapter pointer into the author environment example at source line 865; invoked at proof line 5681 of the counted proof. It is retained only as a navigation pointer; every load-bearing step is inside the delivered span.\par\medskip
\begin{lemma}
\label{lemma-duality-bootstrap}
Let $(S, \omega_S^\bullet)$ be as in Situation \ref{situation-dualizing}.
With $f^!_{new}$ and $\omega_X^\bullet$ defined for all (morphisms of)
schemes of finite type over $S$ as above:
\begin{enumerate}
\item the functors $f^!_{new}$ and the arrows
$(g \circ f)^!_{new} \to f^!_{new} \circ g^!_{new}$
turn $D_{\textit{Coh}}^+$ into a pseudo functor from the category of
schemes of finite type over $S$ into the $2$-category of categories,
\item $\omega_X^\bullet = (X \to S)^!_{new} \omega_S^\bullet$,
\item the functor $D_X$
defines an involution of $D_{\textit{Coh}}(\mathcal{O}_X)$
switching $D_{\textit{Coh}}^+(\mathcal{O}_X)$ and
$D_{\textit{Coh}}^-(\mathcal{O}_X)$ and fixing
$D_{\textit{Coh}}^b(\mathcal{O}_X)$,
\item $\omega_X^\bullet = f^!_{new}\omega_Y^\bullet$ for
$f : X \to Y$ a morphism of finite type schemes over $S$,
\item $f^!_{new}M = D_X(Lf^*D_Y(M))$ for
$M \in D_{\textit{Coh}}^+(\mathcal{O}_Y)$, and
\item if in addition $f$ is proper, then $f^!_{new}$ is isomorphic
to the restriction of the right adjoint of
$Rf_* : D_\QCoh(\mathcal{O}_X) \to D_\QCoh(\mathcal{O}_Y)$ to
$D_{\textit{Coh}}^+(\mathcal{O}_Y)$ and there is a canonical isomorphism
$$
Rf_*R\SheafHom_{\mathcal{O}_X}(K, f^!_{new}M)
\to
R\SheafHom_{\mathcal{O}_Y}(Rf_*K, M)
$$
for $K \in D^-_{\textit{Coh}}(\mathcal{O}_X)$ and
$M \in D_{\textit{Coh}}^+(\mathcal{O}_Y)$, and
$$
Rf_*R\SheafHom_{\mathcal{O}_X}(K, \omega_X^\bullet) =
R\SheafHom_{\mathcal{O}_Y}(Rf_*K, \omega_Y^\bullet)
$$
for $K \in D^-_{\textit{Coh}}(\mathcal{O}_X)$ and
\end{enumerate}
If $X$ is separated over $S$, then
$\omega_X^\bullet$ is canonically isomorphic to
$(X \to S)^!\omega_S^\bullet$ and
if $f$ is a morphism between schemes separated
over $S$, then there is a canonical isomorphism\footnote{We haven't
checked that these are compatible with the isomorphisms
$(g \circ f)^! \to f^! \circ g^!$ and
$(g \circ f)^!_{new} \to f^!_{new} \circ g^!_{new}$. We will do this
here if we need this later.}
$f_{new}^!K = f^!K$ for $K$ in $D_{\textit{Coh}}^+$.
\end{lemma}

\begin{proof}
Let $f : X \to Y$, $g : Y \to Z$, $h : Z \to T$ be morphisms of schemes
of finite type over $S$. We have to show that
$$
\xymatrix{
(h \circ g \circ f)^!_{new} \ar[r] \ar[d] &
f^!_{new} \circ (h \circ g)^!_{new} \ar[d] \\
(g \circ f)^!_{new} \circ h^!_{new} \ar[r] &
f^!_{new} \circ g^!_{new} \circ h^!_{new}
}
$$
is commutative. Let $\eta_Y : \text{id} \to D_Y^2$
and $\eta_Z : \text{id} \to D_Z^2$ be the canonical isomorphisms
of Lemma \ref{lemma-dualizing-schemes}. Then, using
Categories, Lemma \ref{categories-lemma-properties-2-cat-cats},
a computation (omitted) shows that both arrows
$(h \circ g \circ f)^!_{new} \to f^!_{new} \circ g^!_{new} \circ h^!_{new}$
are given by
$$
1 \star \eta_Y \star 1 \star \eta_Z \star 1 :
D_X \circ Lf^* \circ Lg^* \circ Lh^* \circ D_T
\longrightarrow
D_X \circ Lf^* \circ D_Y^2 \circ Lg^* \circ D_Z^2 \circ Lh^* \circ D_T
$$
This proves (1). Part (2) is immediate from the definition of
$(X \to S)^!_{new}$ and the fact that $D_S(\omega_S^\bullet) = \mathcal{O}_S$.
Part (3) is Lemma \ref{lemma-dualizing-schemes}.
Part (4) follows by the same argument as part (2).
Part (5) is the definition of $f^!_{new}$.

\medskip\noindent
Proof of (6). Let $a$ be the
right adjoint of Lemma \ref{lemma-twisted-inverse-image} for the
proper morphism $f : X \to Y$ of schemes of finite type over $S$.
The issue is that we do not know $X$ or $Y$ is
separated over $S$ (and in general this won't be true)
hence we cannot immediately apply
Lemma \ref{lemma-shriek-via-duality} to $f$ over $S$.
To get around this we use the canonical identification
$\omega_X^\bullet = a(\omega_Y^\bullet)$ of
Lemma \ref{lemma-proper-map-good-dualizing-complex}.
Hence $f^!_{new}$ is the restriction of $a$ to
$D_{\textit{Coh}}^+(\mathcal{O}_Y)$ by Lemma \ref{lemma-shriek-via-duality}
applied to $f : X \to Y$ over the base scheme $Y$!
The displayed equalities hold by
Example \ref{example-iso-on-RSheafHom-noetherian}.

\medskip\noindent
The final assertions follow from the construction of normalized
dualizing complexes and the already used Lemma \ref{lemma-shriek-via-duality}.
\end{proof}

\end{document}
